% =====================================================================
%  Variante del diagrama de arquitectura para el paper CoNaIISI 2026.
%  Derivada de ../diagrama_arquitectura.tex (source of truth de la tesina).
%  Diferencias: sin leyenda ni tabla de referencia de puertos (el caption del
%  paper cubre la leyenda y el texto no usa los números de puerto), sin
%  numeración de puertos, cajas de filtro más anchas y fuentes un escalón
%  más grandes, para que sea legible impreso a ancho de página.
%  Compilar: lualatex diagrama_paper.tex  ->  diagrama_paper.pdf (vectorial).
% =====================================================================
\documentclass[tikz,border=5pt]{standalone}
\usepackage{fontspec}
% Times New Roman real desde las fuentes de Windows (WSL); respaldo: TeX Gyre Termes.
\setmainfont{Times New Roman}[Path=/mnt/c/Windows/Fonts/, Extension=.ttf,
  UprightFont=times, BoldFont=timesbd, ItalicFont=timesi, BoldItalicFont=timesbi]
\usepackage[spanish,es-tabla,es-noindentfirst]{babel}
\usetikzlibrary{shapes.geometric,positioning,arrows.meta,calc,fit,backgrounds}

\begin{document}
\begin{tikzpicture}[
  font=\normalsize,
  filtro/.style={draw, thick, rectangle, rounded corners=2pt,
    minimum width=2.5cm, minimum height=1.6cm, align=center},
  exterior/.style={draw, thick, cylinder, shape border rotate=90, aspect=0.25,
    minimum width=2.2cm, minimum height=1.5cm, align=center, font=\normalsize\itshape},
  repos/.style={draw, line width=2pt, rectangle, rounded corners=3pt, align=center, fill=black!4},
  puerto/.style={fill=black, draw=black, rectangle, minimum size=5pt, inner sep=0pt},
  tubo/.style={thick},
  conexion/.style={thick, dash pattern=on 4pt off 3pt},
  sint/.style={{Circle[length=2.2mm]}-{Latex[length=2.8mm,width=2.2mm]}, thick},
  memoria/.style={{Latex[length=3.4mm,width=3mm]}-{Latex[length=3.4mm,width=3mm]}, line width=1.9pt},
  sintbox/.style={draw, thick, rectangle, rounded corners=2pt, align=center,
    font=\footnotesize, fill=black!3, minimum height=0.9cm, text width=2.8cm, inner sep=3pt},
  compuesto/.style={draw, thick, dash pattern=on 3pt off 2pt, rounded corners=6pt},
  resp/.style={text=red, font=\small\bfseries, inner sep=1pt},
]

% ---------- Filtros (banda media, y=4.5; paso 3.1 cm) ----------
\node[filtro] (extr)  at (3,4.5)    {Extracción};
\node[filtro] (enriq) at (6.1,4.5)  {Enriquecimiento\\semántico};
\node[filtro] (comp)  at (9.2,4.5)  {Composición\\de contexto};
\node[filtro] (invoc) at (12.3,4.5) {Invocación\\generativa};
\node[filtro] (deco)  at (15.4,4.5) {Decodificación};
\node[filtro] (verif) at (18.5,4.5) {Verificación};
\node[filtro] (norm)  at (21.6,4.5) {Normalización};

% ---------- Fuente y sumidero ----------
\node[exterior] (portales) at (0.0,4.5) {Portales\\digitales};
\node[exterior] (presenta) at (24.6,4.5) {Capa de\\presentación\\\footnotesize(tableros, reportes)};

% ---------- Puertos (cuadraditos en bordes O/E) ----------
\foreach \f in {extr,enriq,comp,invoc,deco,verif,norm}{
  \node[puerto] (\f-in)  at (\f.west) {};
  \node[puerto] (\f-out) at (\f.east) {};
}

% ---------- Tubos tipados entre filtros ----------
\draw[tubo] (extr-out)  -- (enriq-in);
\draw[tubo, draw=black!45] (enriq-out) -- (comp-in);
\draw[tubo] (comp-out)  -- (invoc-in);
\draw[tubo] (invoc-out) -- (deco-in);
\draw[tubo] (deco-out)  -- (verif-in);
\draw[tubo] (verif-out) -- (norm-in);

% ---------- Conexiones al entorno ----------
\draw[conexion] (portales.east) -- (extr-in);
\draw[conexion] (norm-out) -- (presenta.west);

% ---------- Sintonizadores ----------
\node[sintbox] (sint-extr)  at (3,7.4)    {Selección de portales,\\reglas y prioridad};
\node[sintbox] (sint-comp)  at (9.2,7.4)  {Template, variables\\y parámetros de recuperación};
\node[sintbox] (sint-verif) at (18.5,7.4) {Criterio de validez\\y ejecutor};
\node[sintbox] (sint-norm)  at (21.6,7.4) {Criterios de\\consolidación};
\draw[sint] (sint-extr.south)  -- (extr.north);
\draw[sint] (sint-comp.south)  -- (comp.north);
\draw[sint] (sint-verif.south) -- (verif.north);
\draw[sint] (sint-norm.south)  -- (norm.north);

% ---------- Filtro compuesto: Recuperación aumentada ----------
\begin{scope}[on background layer]
  \node[compuesto, fit=(enriq)(comp), inner xsep=7pt, inner ysep=12pt] (rag) {};
\end{scope}
\node[font=\normalsize\itshape, anchor=south west] at (rag.north west) {Recuperación aumentada};
\node[puerto, fill=black!40, draw=black!40] at (enriq.east) {};
\node[puerto, fill=black!40, draw=black!40] at (comp.west) {};

% ---------- Repositorio común ----------
\node[repos, minimum width=21.1cm, minimum height=1.2cm] (repo) at (12.3,0.3)
  {Repositorio común --- PostgreSQL / ORM de Odoo};
\foreach \f in {extr,enriq,comp,invoc,deco,verif,norm}{
  \draw[memoria] (\f.south) -- (\f.south |- repo.north);
}

% ---------- Marcadores de responsabilidades [A]--[F] ----------
\node[resp, anchor=west] at ($(extr.south)  +(0.12,-0.74)$) {[A]\,[F]};
\node[resp, anchor=west] at ($(enriq.south) +(0.12,-0.74)$) {[C]};
\node[resp, anchor=west] at ($(comp.south)  +(0.12,-0.74)$) {[A]\,[B]};
\node[resp, anchor=west] at ($(invoc.south) +(0.12,-0.74)$) {[C]};
\node[resp, anchor=west] at ($(deco.south)  +(0.12,-0.74)$) {[B]\,[C]};
\node[resp, anchor=west] at ($(verif.south) +(0.12,-0.74)$) {[A]\,[B]\,[D]};
\node[resp, anchor=west] at ($(norm.south)  +(0.12,-0.74)$) {[A]\,[B]\,[D]};
\node[resp, anchor=north] at ($(portales.south)+(0,-0.1)$) {[F]};
\node[resp, anchor=east]  at ($(repo.east)+(-0.25,0)$) {[E]};

\end{tikzpicture}
\end{document}
